%% kapa: exam=1 round=36 year=2025 period=1 form=A subject=civil_law subject_ko=민법 seq=30 number=30
%% source: Q-Net 감정평가사 기출문제 (cst00309) artlSeq=5231525 file=raw/1008/5231525/2025년 제36회 감정평가사 1차시험 1교시 문제지 원본.pdf
%% conversion: kordoc 4.24.51 -> tools/gen1.py (auto); review=proofread-sonnet; flags=none
\begin{question}{30}
소유물반환청구권에 관한 설명으로 옳은 것을 모두 고른 것은? (다툼이 있으면 판례에 따름)
\begin{framedtext}
ㄱ. 타인의 소유물을 불법으로 점유하였던 자라도 더 이상 현실적으로 점유를 하고 있지 않은 이상 그를 상대로 한 소유물반환청구는 부당하다. \par
ㄴ. 타인의 토지에 무단으로 건물을 신축하여 소유하는 자에 대하여 토지소유자는 그 건물에서 퇴거할 것을 청구할 수 있다. \par
ㄷ. 토지소유자는 토지에 대한 점유취득시효 완성자에 대하여 불법점유를 이유로 토지의 반환을 청구할 수 있다.
\end{framedtext}
\begin{choices}
  \item ㄱ
  \item ㄴ
  \item ㄷ
  \item ㄱ, ㄴ
  \item ㄴ, ㄷ
\end{choices}
\end{question}
\answer{1}
